% SPDX-License-Identifier: CC-BY-SA-4.0
% Author: Matthieu Perrin
% Part: 
% Section: 
% Sub-section: 
% Frame: 

\begingroup

\begin{frame}[fragile]{Lemme 1 : $A$ possède une configuration critique}

  \begin{exampleblock}{Démonstration}
    \begin{itemize}
    \item Supposons (absurde), que $A$ ne possède pas de configuration critique.
    \item Construction par récurrence :\\
      une exécution {\color{exampleColor}infinie} qui reste dans des {\color{exampleColor}configurations bivalentes}
      \begin{description}
      \item[Initialisation] La configuration initiale est bivalente
        \begin{itemize}
        \item Si $T_0$ crashe au début, $T_1$ décide $\alert{\bullet}$
        \item Si $T_1$ crashe au début, $T_0$ décide $\structure{\bullet}$
        \end{itemize}
      \item[Hérédité] Supposons qu'on a construit l'exécution $C_1, ..., C_k$
        \begin{itemize}
        \item $C_k$ est bivalente
        \item $C_k$ n'est pas critique
        \item $C_k$ possède un fils bivalent $C_{k+1}$ 
        \end{itemize}
      \end{description}
    \item $A$ est wait-free mais possède une exécution infinie : {\color{exampleColor}absurde !}
    \end{itemize}
  \end{exampleblock}

\end{frame}

\endgroup
\endinput
